% !TeX root = ../main.tex

\chapter{Introduction}
\label{chap:intro}

Distinguishing non-orthgonal quantum states is a very important topic in quantum
information since it is the foundation of the first quantum cryptography protocols,
such as BB84 \cite{bb84} and SARG04 \cite{RN492}.

\begin{quote}
    When information is encoded in non-orthogonal quantum
    states, such as single photons with polarization directions
    0, 45, 90, and 135 degrees, one obtains a communications channel whose
    transmissions in principle cannot be read or copied reliably by an
    eavesdropper ignorant of certain key information used in forming the transmission.
\end{quote}

Most of the quantum states involved in these protocols
have easier structures, such as orthogonal states with basis change
or trine states, since the measurement for these sets of states
are readily available or easier to design than other states.
Quantum state discrimination in those protocols are more like
an attempt to crack the eavesdropped messages and a proof that
such protocols are safe.

In a more general scenario, people try to discriminate a set of
quantum states that are non-orthogonal and linearly independent.
Defining and finding a "good" measurement for these states became
non-trivial. Two important discrimination schemes were proposed,
minimum error state discrimination and unambiguous state discrimination.
In the protocol B92, USD plays an important role as a step for Bob.

Several works have discussed the condition where noise is present before or
during the discrimination process.

\section{Our motivation}
% This paragraph is produced with the help of ChatGPT 4o.
Quantum transduction
\cite{benbrubakerBuildQuantum,ibmQuantumTransduction}
has emerged as a key interface between distinct physical platforms,
particularly superconducting processors and optical communication channels
\cite{RN602,RN648}.
While not the primary focus of this work, it illustrates the broader effort to
integrate quantum computing and communication, motivating our exploration of
how optical communication scenarios can be expressed in circuit-level terms.

The first question we asked is how to construct quantum circuits for quantum state
discrimination. It turns out that it is not trivial, but it is also not too difficult.
By Naimark's dilation theorem (See Sec \ref{sec:naimark}), it is possible to
convert any POVM with finite elements to projective measurements or a unitary
operation.
Then by applying the state-of-the-art quantum circuit synthesis tools, we can
verify that the quantum circuits successfully realize the functionality of
input state discrimination.

After the construction of the quantum circuits, we asked how noise will affect
the discrimination results.
We reviewed the discussions of quantum state discrimination of mixed states.
We were curious whether there are better definitions of good measurements than
existing methods for unambiguous state discrimination.
We proposed a series of constraints, and we will show our preliminary results.
This is under the assumption that the quantum circuit implementation is
noiseless. As for noisy quantum circuits, we will in turn assume the input
quantum states are noiseless.

% A good reason why we discuss multi-qubit USD. \cite{RN485}

\section{Our contributions}
\begin{itemize}
    \item New definitions for error-tolerant quantum state discrimination
    \item A quantum circuit synthesis flow for quantum state discrimination
\end{itemize}

\section{Thesis organization}
This thesis is structured as follows.
Chapter \ref{chap:basics_qc} introduces the fundamentals of quantum computation,
providing the necessary background for circuit-based implementations.
Chapter \ref{chap:basics_qsd} reviews quantum state discrimination, covering
key concepts like MED and UQSD.
Chapter \ref{chap:basics_co} outlines convex optimization techniques used in
our optimization strategies.
In Chapter \ref{chap:error_uqsd}, we develop noise-aware approaches for
quantum state discrimination, addressing challenges in CrossQSD and FitQSD.
Chapter \ref{chap:qc_qsd} presents our quantum circuit synthesis flow for
state discrimination and analyzes the impact of noise on circuits and input states.
Finally, Chapter \ref{chap:summary} summarizes our findings and discusses
future research directions to advance practical quantum state discrimination.
